\section{A positive two-sector transfer model}

To test the vacuum interpretation rather than only Perron optimization, take
\begin{equation}
 T_\varepsilon=
 \begin{pmatrix}
 1&\varepsilon\\
 \varepsilon&2\varepsilon^2
 \end{pmatrix},
 \qquad 0<\varepsilon<2^{-1/2}.
\end{equation}
It is symmetric, entrywise nonnegative, and positive definite because its
determinant is \(\varepsilon^2\).  Declare a bad extension to be a transition
whose next state is sector one.

At \(\varepsilon=0.002430277501841998\), the certificate gives
\begin{align}
 \lambda_0&=1.0000059062836195,\\
 \psi&=(1,\;0.002430291855665421),\\
 V=\psi^{-1}&=(1,\;411.4732136672519).
\end{align}
The Doob row-sum error is below \(2.3\times10^{-16}\), the stationary and
detailed-balance checks pass at the same scale, and
\begin{equation}
 \max_i\frac{(QV)_i}{V_i}=0.002430263148018576.
\end{equation}
The identity~\eqref{eq:inverse-cancel} holds with error below
\(4.4\times10^{-19}\).

Using \(\psi\) instead of \(\psi^{-1}\) in the forward cost inequality does
not yield the cancellation identity; the recorded error is
\(2.4302488\times10^{-3}\).  This does not say that amplitude weighting is
never useful.  It says that the two orientations answer different questions
and cannot be interchanged silently.

\subsection{Finite positivity is not uniformity}

Let a sequence of finite volumes have a rare-sector amplitude
\(\psi_L=e^{-L}\).  Every finite-volume vacuum is positive, but the inverse
cost weight grows as \(e^L\).  At \(L=20\) it exceeds
\(4.85\times10^8\).  Perron positivity alone therefore supplies neither a
volume-uniform Harnack bound nor a uniform root moment.  Those are separate
analytic obligations.
